\documentclass[a4paper,12pt]{article}

\usepackage[T2A]{fontenc}
\usepackage[utf8]{inputenc}
\usepackage[russian]{babel}
\usepackage{amsmath,amssymb,amsthm,mathtools}
\usepackage{helvet}
\renewcommand{\familydefault}{\sfdefault}
\usepackage{icomma}
\usepackage{geometry}
\geometry{left=20mm,right=20mm,top=20mm,bottom=22mm}
\usepackage{microtype}
\usepackage{tikz}
\usetikzlibrary{calc,arrows.meta,positioning,shapes.geometric}
\usepackage{tabularx,booktabs,longtable}
\usepackage{enumitem}
\usepackage{hyperref}
\usepackage{parskip}
\usepackage{fancyhdr}
\usepackage{setspace}
\usepackage{xcolor}

\definecolor{accent}{HTML}{008080}
\definecolor{darkaccent}{HTML}{004D4D}
\definecolor{textgray}{HTML}{333333}
\definecolor{softgray}{HTML}{6B7280}

\hypersetup{hidelinks}

\onehalfspacing
\setlength{\parindent}{0pt}
\setlength{\headheight}{14pt}
\setlength{\emergencystretch}{2em}

\setlist[itemize]{
  leftmargin=1.5em,
  itemsep=0.25em,
  topsep=0.35em
}

\setlist[enumerate]{
  leftmargin=1.8em,
  itemsep=0.35em,
  topsep=0.35em
}

\renewcommand{\arraystretch}{1.35}

\pagestyle{fancy}
\fancyhf{}
\fancyhead[L]{\small\color{softgray} Математика, 6 класс}
\fancyhead[R]{\small\color{softgray} Надя Клименко}
\fancyfoot[C]{\small\color{softgray}\thepage}
\renewcommand{\headrulewidth}{0.3pt}
\renewcommand{\footrulewidth}{0pt}

\newcommand{\sectionrule}{%
  \par
  \vspace{-0.35em}
  {\color{accent}\rule{\linewidth}{0.7pt}}
  \vspace{0.55em}
}

\newcommand{\key}[1]{\textbf{\color{darkaccent}#1}}
\newcommand{\term}[1]{\textbf{#1}}

\tikzset{
  flowstart/.style={
    ellipse,
    draw=darkaccent,
    line width=0.9pt,
    minimum width=4.2cm,
    minimum height=1.05cm,
    align=center,
    text width=3.8cm
  },
  flowproc/.style={
    rectangle,
    rounded corners=3pt,
    draw=accent,
    line width=0.9pt,
    minimum height=1.15cm,
    align=center,
    text width=5.5cm,
    inner sep=7pt
  },
  flowcond/.style={
    diamond,
    aspect=2.1,
    draw=darkaccent,
    line width=0.9pt,
    align=center,
    text width=3.8cm,
    inner sep=2pt
  },
  flowarrow/.style={
    -{Stealth[length=3mm,width=2mm]},
    line width=0.9pt,
    draw=textgray
  },
  flowlabel/.style={
    font=\small,
    fill=white,
    inner sep=1.5pt,
    text=textgray
  }
}

\begin{document}

\color{textgray}

% =================================================
% ТИТУЛЬНЫЙ ЛИСТ
% =================================================

\begin{titlepage}
\thispagestyle{empty}

\centering

\vspace*{2.5cm}

{\Large\color{darkaccent}\textbf{Математика, 6 класс}\par}

\vspace{1.3cm}

{\Huge\textbf{Чек-лист}\par}

\vspace{0.55cm}

{\LARGE\color{accent}
\textbf{Средняя скорость, математические модели и проценты}\par}

\vspace{1.7cm}

{\large Надя Клименко\par}

\vspace{0.4cm}

{\large \today\par}

\vfill

{\large\textbf{Лёвин Артём Александрович}\par}

\vspace{0.2cm}

{\normalsize эксклюзивно для Нади Клименко\par}

\vspace{1.0cm}

\begin{tikzpicture}[x=1cm,y=1cm]
\draw[
  accent,
  line width=1.1pt
]
(-7.2,0)
.. controls (-4.8,0.8) and (-2.4,-0.7) .. (0,0)
.. controls (2.4,0.7) and (4.8,-0.8) .. (7.2,0);

\draw[
  darkaccent,
  line width=0.45pt
]
(-6.4,-0.25)
.. controls (-3.8,0.35) and (-1.4,-0.35) .. (1.1,-0.05)
.. controls (3.2,0.22) and (4.8,-0.2) .. (6.4,0.05);
\end{tikzpicture}

\end{titlepage}

% =================================================
% СРЕДНЯЯ СКОРОСТЬ
% =================================================

\section*{1. Средняя скорость}
\sectionrule

\term{Средняя скорость} показывает, с какой постоянной скоростью
нужно было бы двигаться, чтобы пройти \key{тот же весь путь}
за \key{то же всё время}.

\[
\boxed{
v_{\text{ср}}
=
\frac{S_{\text{весь}}}{t_{\text{всё}}}
}
\]

\key{Главная мысль:}
сначала складывают все пройденные расстояния,
затем складывают всё время движения.
Только после этого общий путь делят на общее время.

\textbf{Пример.}
Первые $3$ ч тело двигалось со скоростью $65$ км/ч,
следующие $2$ ч --- со скоростью $88$ км/ч.

\[
S_1=65\cdot3=195\text{ км},
\qquad
S_2=88\cdot2=176\text{ км}.
\]

\[
S_{\text{весь}}
=
195+176
=
371\text{ км}.
\]

\[
t_{\text{всё}}
=
3+2
=
5\text{ ч}.
\]

\[
v_{\text{ср}}
=
\frac{371}{5}
=
74{,}2\text{ км/ч}.
\]

\textbf{Смысл результата.}
Если бы тело всё время двигалось с постоянной скоростью
$74{,}2$ км/ч, оно прошло бы те же $371$ км за те же $5$ ч.

\textbf{Важно.}

\begin{itemize}
  \item Нельзя автоматически вычислять
  \[
  \frac{v_1+v_2}{2}.
  \]
  Такой приём работает лишь в специальных случаях.

  \item Перед вычислением все единицы времени должны быть согласованы.

  \item Например:
  \[
  30\text{ мин}=0{,}5\text{ ч},
  \qquad
  12\text{ мин}=\frac{12}{60}\text{ ч}=0{,}2\text{ ч}.
  \]

  \item Средняя скорость относится ко всему движению целиком.
\end{itemize}

% =================================================
% БЛОК-СХЕМА: СРЕДНЯЯ СКОРОСТЬ
% =================================================

\clearpage
\thispagestyle{fancy}

\section*{Алгоритм: как найти среднюю скорость}

\vspace{0.8cm}

\begin{center}
\begin{tikzpicture}[node distance=11mm and 18mm]

\node[flowstart] (start)
{Начало};

\node[flowproc,below=of start] (read)
{Выпишите все участки пути и соответствующее им время};

\node[flowproc,below=of read] (units)
{Приведите время к одинаковым единицам};

\node[flowproc,below=of units] (sumS)
{Найдите общий путь:
$S_{\text{весь}}=S_1+S_2+\dots$};

\node[flowproc,below=of sumS] (sumT)
{Найдите всё время:
$t_{\text{всё}}=t_1+t_2+\dots$};

\node[flowproc,below=of sumT] (calc)
{Вычислите
$v_{\text{ср}}=
\dfrac{S_{\text{весь}}}{t_{\text{всё}}}$};

\node[flowproc,below=of calc] (check)
{Проверьте единицы измерения и смысл результата};

\node[flowstart,below=of check] (end)
{Ответ};

\draw[flowarrow] (start) -- (read);
\draw[flowarrow] (read) -- (units);
\draw[flowarrow] (units) -- (sumS);
\draw[flowarrow] (sumS) -- (sumT);
\draw[flowarrow] (sumT) -- (calc);
\draw[flowarrow] (calc) -- (check);
\draw[flowarrow] (check) -- (end);

\end{tikzpicture}
\end{center}

\clearpage

% =================================================
% СРЕДНЕЕ АРИФМЕТИЧЕСКОЕ И МОДЕЛИ
% =================================================

\section*{2. Среднее арифметическое и математическая модель}
\sectionrule

Для двух чисел $x$ и $y$ среднее арифметическое равно

\[
\boxed{
\frac{x+y}{2}
}
\]

Если одно число в $3$ раза больше другого,
удобно обозначить меньшее число через $x$,
а большее --- через $3x$.

\textbf{Пример.}
Среднее арифметическое двух чисел равно $36$,
одно число в $3$ раза больше другого.

\[
\frac{x+3x}{2}=36.
\]

Приводим подобные:

\[
x+3x=(1+3)x=4x.
\]

Получаем

\[
\frac{4x}{2}=36,
\]

\[
2x=36,
\]

\[
x=18.
\]

Второе число:

\[
3x=54.
\]

Ответ:

\[
18,\quad 54.
\]

\subsection*{Приведение подобных}

Если буквенная часть одинаковая,
складывают или вычитают коэффициенты:

\[
x+3x=4x,
\]

\[
7a-2a=5a,
\]

\[
4m+m=5m.
\]

\key{Важно:}
складываются коэффициенты при одинаковой буквенной части.

% =================================================

\subsection*{Последовательные натуральные числа}

Три последовательных натуральных числа удобно записывать так:

\[
a,\qquad a+1,\qquad a+2.
\]

\textbf{Пример.}
Сумма трёх последовательных натуральных чисел равна $114$.

Составляем уравнение:

\[
a+(a+1)+(a+2)=114.
\]

Приводим подобные:

\[
3a+3=114.
\]

\[
3a=111.
\]

\[
a=37.
\]

Следовательно,

\[
37,\qquad38,\qquad39.
\]

Проверка:

\[
37+38+39=114.
\]

\subsection*{Типичные ошибки}

\begin{itemize}
  \item Записать
  \[
  x+3x=2x.
  \]
  Верно:
  \[
  x+3x=4x.
  \]

  \item Для трёх последовательных чисел записать
  $a$, $a$, $a$.

  \item Найти значение переменной и забыть восстановить остальные числа.

  \item Составить уравнение без связи с текстом задачи.
\end{itemize}

% =================================================
% ПРОЦЕНТ
% =================================================

\section*{3. Что такое процент}
\sectionrule

\term{Один процент} --- сотая доля величины:

\[
\boxed{
1\%=\frac{1}{100}=0{,}01
}
\]

Поэтому $p\%$ от числа $a$ можно найти двумя
равносильными способами:

\[
\boxed{
a\cdot\frac{p}{100}
}
\]

или

\[
\boxed{
\frac{a}{100}\cdot p
}.
\]

\textbf{Первый путь: перевод процента в дробь.}

Например,

\[
50\%
=
\frac{50}{100}
=
\frac12.
\]

Поэтому

\[
50\%\text{ от }200
=
200\cdot\frac12
=
100.
\]

\textbf{Второй путь: сначала найти $1\%$.}

\[
1\%\text{ от }200
=
\frac{200}{100}
=
2.
\]

Тогда

\[
50\%\text{ от }200
=
2\cdot50
=
100.
\]

Оба метода приводят к одному результату.

% =================================================

\subsection*{Полезные быстрые соответствия}

\begin{center}

\begin{tabularx}{\linewidth}{
@{}
>{\centering\arraybackslash}X
>{\centering\arraybackslash}X
>{\centering\arraybackslash}X
@{}
}
\toprule
\textbf{Процент}
&
\textbf{Доля}
&
\textbf{Быстрое действие}
\\
\midrule

$50\%$
&
$\dfrac12$
&
разделить на $2$
\\

$25\%$
&
$\dfrac14$
&
разделить на $4$
\\

$20\%$
&
$\dfrac15$
&
разделить на $5$
\\

$10\%$
&
$\dfrac1{10}$
&
разделить на $10$
\\

$5\%$
&
$\dfrac1{20}$
&
разделить на $20$
\\

$1\%$
&
$\dfrac1{100}$
&
разделить на $100$
\\

\bottomrule
\end{tabularx}

\end{center}

Эти соответствия полезны для устного счёта и самопроверки.
При первичном изучении темы надёжно сначала видеть связь
процента с дробью.

% =================================================
% СКИДКИ
% =================================================

\section*{4. Скидка как задача на проценты}
\sectionrule

Если товар стоит $1000$ руб.,
а скидка составляет $5\%$,
сначала определяют размер скидки.

\[
5\%\text{ от }1000
=
1000\cdot\frac{5}{100}
=
50\text{ руб.}
\]

После этого находят итоговую цену:

\[
1000-50=950\text{ руб.}
\]

\key{Различайте два вопроса.}

\begin{itemize}
  \item
  «Сколько составляет скидка?» ---
  требуется найти процент от исходной цены.

  \item
  «Сколько нужно заплатить?» ---
  требуется вычесть найденную скидку
  из исходной цены.
\end{itemize}

\textbf{Ещё один пример.}

Шапка стоит $2000$ руб.,
скидка равна $10\%$.

\[
10\%\text{ от }2000
=
200.
\]

\[
2000-200=1800\text{ руб.}
\]

Шарф стоит $5000$ руб.
При той же скидке:

\[
10\%\text{ от }5000=500,
\]

\[
5000-500=4500\text{ руб.}
\]

Если скидка действует на оба товара,
можно сначала сложить стоимости:

\[
2000+5000=7000.
\]

Затем вычислить общую скидку:

\[
10\%\text{ от }7000=700.
\]

Итоговая стоимость:

\[
7000-700=6300\text{ руб.}
\]

% =================================================
% ТРИ ТИПА ЗАДАЧ
% =================================================

\section*{5. Три основных типа задач на проценты}
\sectionrule

Обозначим:

\begin{itemize}
  \item
  $a$ --- целое, число, \emph{от которого} считают процент;

  \item
  $v$ --- часть, число, \emph{для которого} определяют процент;

  \item
  $p$ --- количество процентов.
\end{itemize}

% -------------------------------------------------

\subsection*{Тип 1. Найти $p\%$ от числа $a$}

\[
\boxed{
v=a\cdot\frac{p}{100}
}
\]

\textbf{Пример.}
В классе $30$ учеников.
$10\%$ изучают немецкий язык.

\[
30\cdot\frac{10}{100}=3.
\]

Значит, немецкий язык изучают $3$ ученика.

Не изучают немецкий язык:

\[
30-3=27.
\]

% -------------------------------------------------

\subsection*{Тип 2. Найти целое $a$, если известны часть $v$ и процент $p$}

\[
\boxed{
a=\frac{100v}{p}
}
\]

\textbf{Пример.}
$3$ ученика составляют $10\%$ всех учеников класса.

\[
a
=
\frac{100\cdot3}{10}
=
30.
\]

Если те же $3$ ученика составляют $5\%$,
получаем

\[
a
=
\frac{100\cdot3}{5}
=
60.
\]

% -------------------------------------------------

\subsection*{Тип 3. Найти, сколько процентов число $v$ составляет от числа $a$}

Сначала определяют долю части в целом:

\[
\frac{v}{a}.
\]

Затем переводят эту долю в проценты:

\[
\boxed{
p=\frac{v}{a}\cdot100\%
}
\]

\textbf{Пример.}
$3$ ученика из $30$ изучают немецкий язык.

\[
p
=
\frac{3}{30}\cdot100\%
=
10\%.
\]

\key{Ключевой смысл:}

\[
\frac{\text{часть}}{\text{целое}}
\]

показывает долю части в целом.

% =================================================
% БЛОК-СХЕМА: ПРОЦЕНТЫ
% =================================================

\clearpage
\thispagestyle{fancy}

\section*{Алгоритм: как определить тип задачи на проценты}

\vspace{0.8cm}

\begin{center}

\begin{tikzpicture}[
scale=0.9,
transform shape,
node distance=14mm and 12mm
]

\node[flowstart] (start)
{Прочитайте условие};

\node[flowproc,below=of start] (mark)
{Определите: где целое $a$, где часть $v$, где процент $p$};

\node[flowcond,below=18mm of mark] (question)
{Что требуется найти?};

\node[
  flowproc,
  below left=20mm and 31mm of question,
  text width=4.1cm
] (findv)
{Часть $v$\\[2pt]
$v=a\cdot\dfrac{p}{100}$};

\node[
  flowproc,
  below=20mm of question,
  text width=4.1cm
] (finda)
{Целое $a$\\[2pt]
$a=\dfrac{100v}{p}$};

\node[
  flowproc,
  below right=20mm and 31mm of question,
  text width=4.1cm
] (findp)
{Процент $p$\\[2pt]
$p=\dfrac{v}{a}\cdot100\%$};

\node[
  flowproc,
  below=28mm of finda,
  text width=5.8cm
] (check)
{Проверьте смысл:
часть и целое выбраны верно,
единицы согласованы};

\node[flowstart,below=of check] (end)
{Ответ};

\draw[flowarrow]
(start) -- (mark);

\draw[flowarrow]
(mark) -- (question);

\draw[flowarrow]
(question.west)
-| node[flowlabel,pos=0.25]{найти $v$}
(findv.north);

\draw[flowarrow]
(question.south)
-- node[flowlabel,right]{найти $a$}
(finda.north);

\draw[flowarrow]
(question.east)
-| node[flowlabel,pos=0.25]{найти $p$}
(findp.north);

\draw[flowarrow]
(findv.south)
|- (check.west);

\draw[flowarrow]
(finda.south)
-- (check.north);

\draw[flowarrow]
(findp.south)
|- (check.east);

\draw[flowarrow]
(check) -- (end);

\end{tikzpicture}

\end{center}

\clearpage

% =================================================
% PVA
% =================================================

\section*{6. Формула PVA как формула-связка}
\sectionrule

Все три величины связаны одной формулой:

\[
\boxed{
p=\frac{v}{a}\cdot100\%
}
\]

где

\[
a \text{ --- целое},
\qquad
v \text{ --- часть},
\qquad
p \text{ --- процент}.
\]

Из неё можно получить остальные формы:

\[
\boxed{
v=\frac{ap}{100}
}
\]

и

\[
\boxed{
a=\frac{100v}{p}
}.
\]

Формула PVA удобна как универсальная связь между тремя величинами
и как средство самопроверки.

На первых этапах полезно сначала словами определить,
что является целым, частью и процентом,
затем подставлять числа.

\subsection*{Дробные ответы допустимы}

Например,

\[
15\%\text{ от }50
=
50\cdot\frac{15}{100}
=
7{,}5.
\]

Дробный результат в задаче на проценты
сам по себе не является признаком ошибки.

% =================================================
% ТИПИЧНЫЕ ОШИБКИ
% =================================================

\section*{7. Типичные ошибки в задачах на проценты}
\sectionrule

\begin{itemize}

  \item
  Перепутать часть и целое.

  Например, если $3$ ученика из $30$ изучают язык,
  правильная доля:

  \[
  \frac{3}{30},
  \]

  поскольку $3$ --- часть, $30$ --- целое.

  \item
  Забыть умножить найденную долю на $100\%$,
  когда требуется ответ именно в процентах.

  \item
  В задаче со скидкой найти размер скидки,
  хотя требуется итоговая цена.

  \item
  При поиске целого умножить известную часть на
  $\dfrac{p}{100}$ ещё раз.

  \item
  Механически применять формулу,
  предварительно не разобрав смысл величин.

  \item
  Считать, что процент обязательно должен быть целым числом.

\end{itemize}

% =================================================
% САМОПРОВЕРКА
% =================================================

\section*{8. Самопроверка перед ответом}
\sectionrule

Перед завершением решения проверьте:

\begin{enumerate}

  \item
  Что является целым?

  \item
  Что является частью?

  \item
  Известен ли процент?

  \item
  Если я ищу процент от числа,
  использовал ли я
  \[
  v=a\cdot\frac{p}{100}?
  \]

  \item
  Если я ищу целое по известной части,
  разумно ли, что при $p<100\%$ целое больше части?

  \item
  Если я ищу процент,
  разделил ли я часть на целое?

  \item
  В задаче со скидкой требуется размер скидки
  или итоговая стоимость?

  \item
  В задаче на среднюю скорость использованы
  весь путь и всё время?

  \item
  Все единицы времени приведены к одинаковому виду?

  \item
  В алгебраической модели правильно ли приведены подобные слагаемые?

\end{enumerate}

% =================================================
% ТРЕНИРОВОЧНЫЙ БЛОК
% =================================================

\clearpage

\section*{Тренировочный блок}
\sectionrule

\textbf{Решайте письменно.}
В задачах на проценты сначала определите,
что является целым, частью и процентом.

\subsection*{Средний уровень}

\begin{enumerate}[label=\arabic*)]

  \item
  Турист прошёл $9$ км за $2$ ч,
  а затем ещё $12$ км за $3$ ч.
  Найдите среднюю скорость на всём пути.

  \item
  Среднее арифметическое двух чисел равно $40$.
  Одно число в $3$ раза больше другого.
  Найдите оба числа.

  \item
  Найдите $15\%$ от числа $240$.

\end{enumerate}

\subsection*{Выше среднего}

\begin{enumerate}[label=\arabic*),start=4]

  \item
  Товар стоит $1500$ руб.
  Скидка составляет $8\%$.
  Сколько рублей нужно заплатить после скидки?

  \item
  Число $18$ составляет $12\%$ некоторого числа.
  Найдите это число.

  \item
  Сколько процентов составляет число $27$ от числа $45$?

  \item
  Сумма трёх последовательных натуральных чисел равна $156$.
  Найдите эти числа.

  \item
  Велосипедист проехал $2{,}4$ км за $12$ мин,
  затем $3{,}6$ км за $18$ мин
  и ещё $1{,}5$ км за $10$ мин.
  Найдите среднюю скорость на всём пути в км/ч.

  \item
  В классе $32$ ученика.
  $25\%$ занимаются в математическом кружке.
  Сколько учеников занимается в кружке?
  Сколько учеников в нём не занимается?

\end{enumerate}

\subsection*{Сложная задача}

\begin{enumerate}[label=\arabic*),start=10]

  \item
  После скидки $15\%$ куртка стала стоить $2720$ руб.
  Найдите первоначальную цену куртки
  и размер скидки в рублях.

\end{enumerate}

% =================================================
% ОТВЕТЫ
% =================================================

\clearpage

\section*{Ответы к тренировочному блоку}
\sectionrule

\begin{table}[ht]
\centering
\caption{Краткие ответы}

\begin{tabularx}{\linewidth}{
@{}
>{\centering\arraybackslash}p{1.2cm}
X
@{}
}

\toprule

\textbf{№}
&
\textbf{Ответ}
\\

\midrule

1
&
$4{,}2$ км/ч
\\

2
&
$20$ и $60$
\\

3
&
$36$
\\

4
&
$1380$ руб.
\\

5
&
$150$
\\

6
&
$60\%$
\\

7
&
$51$, $52$, $53$
\\

8
&
$11{,}25$ км/ч
\\

9
&
$8$ учеников занимаются,
$24$ ученика не занимаются
\\

10
&
первоначальная цена $3200$ руб.;
скидка $480$ руб.
\\

\bottomrule

\end{tabularx}
\end{table}

\vfill

{\small\color{softgray}
Ключ к теме:
весь путь / всё время;
часть / целое;
$1\%=\dfrac1{100}$;
сначала смысл, затем формула.
}

\end{document}